% Part: history
% Chapter: biographies 
% Section: alfred-tarski

\documentclass[../../../include/open-logic-section]{subfiles}

\begin{document}

\olfileid{his}{bio}{tar}

\olsection{الفرېډ تارسکي}

\olphoto{tarski-alfred}{الفرېډ تارسکي}

الفرېډ تارسکي د 1901 د جنورۍ پر 14مه
د پولېنډ په وارسا کښې زېږېدلے و؛
هغه وخت دا د روسيې د سترواکۍ برخه
وه۔ تارسکي چې «ناپوليوني» بلل شوے
دے، پُرجوش، ډېر غږېدونکے او توند
و۔ توان ئې ډېر ځله په درسونو کښې
هم ښکاره کېده؛ يو ځل ئې د درس
پر وخت د سګرټ په غورځولو د کثافاتو
ټوکرۍ اور واخيست، او بيا په هغه
ودانۍ کښې له درس ورکولو منع شو۔

تارسکي له کم عمره د پوهې تنده لرله۔
که څه هم وروسته به ئې زده کوونکو
ته ويل چې منطق ئې ځکه ولوست چې
دا يوازينے مضمون و چې B درجه ئې
پکښې اخيستې وه، د ثانوي ښوونځي
اسناد ئې ښيي چې په ټولو مضامينو،
آن په منطق کښې، ئې A درجې اخيستې
وې۔ له 1918 څخه تر 1924 پورې
ئې د وارسا په پوهنتون کښې زده
کړه وکړه۔ لومړۍ اراده ئې د
بيولوجۍ لوستل وو، خو په رياضياتو،
فلسفه او منطق کښې ئې شوق پيدا
شو، ځکه پوهنتون د منطق او فلسفې
د وارسا مکتب مرکز و۔ تارسکي په
1924 کښې د ستانيسواف لېشنيېفسکي
تر لارښوونې لاندې دوکتورا واخيسته۔

په 1939 کښې امريکا ته له کډې وړلو
مخکښې، تارسکي د خپل ځينې تر ټولو
مهم کار بشپړول په وارسا کښې د
ثانوي ښوونځي د ښوونکي په توګه
وکړل۔ د منطقي پايلې او منطقي
رښتياوالي آثار ئې په همدې وخت
کښې وليکل۔ په 1939 کښې تارسکي
د درسونو د يو سفر دپاره امريکا
ته تلے و۔ د هغه د سفر پر وخت
جرمني پر پولېنډ يرغل وکړ، او د
يهودي اصل له امله تارسکي بېرته
نۀ شو ستنېدلے۔ مېرمن او ماشومان
ئې د جګړې تر پايه په پولېنډ کښې
پاتې شول، خو بيا هغوي هم امريکا
ته کډه يووړلے شوه۔ تارسکي په
هارورډ، د نيويارک ښار په کالج،
د پرينسټن د لوړو څېړنو په موسسه
او بالاخره د برکلي د کليفورنيا
په پوهنتون کښې تدريس وکړ۔ هلته
ئې د منطق او د علم د طريقه‌پوهې
ګڼ‌څانګيز پروګرام جوړ کړ۔ تارسکي
د 1983 د اکتوبر پر 26مه د 82
کلونو په عمر وفات شو۔

\begin{reading}
د تارسکي د ژوند په هکله د نورو
معلوماتو دپاره ژوندليک \emph{الفرېډ
تارسکي: ژوند او منطق}
\citep{Feferman2004} وګورئ۔ د منطقي
پايلې او رښتياوالي په هکله ئې
بنسټيز آثار په انګليسي کښې په
\citep{Tarski1983} کښې شته۔ د
تارسکي ټول اصلي آثار په څلور
ټوکي لړۍ \citep{Tarski1981} کښې
راټول شوي دي۔
\end{reading}

\end{document}
